\section{Introduction}
\label{sec:introduction}

At 3:58 on a weekday afternoon, a Manhattan taxi ride ends with a
\$17.70 metered fare. The backseat screen assembles the bill---fare,
fixed fees, congestion surcharge---into a \$21.70 total and offers three
buttons: 20, 25, or 30 percent, computed on that total. Four minutes
later, the identical ride ends differently. The rush-hour surcharge has
switched on, the total reads \$24.20, and every button is higher---the
20 percent button by exactly fifty cents, which is 20 percent of the
\$2.50 the fare schedule just added (Figure~\ref{fig:screens}). Nothing
about the trip, the driver, or the rider's regard for the driver has
changed. A scheduled line item changed, and the suggestion arithmetic
carried it into the tip.

\subsection{The mechanism on the screen}
Figure~\ref{fig:feestack} shows what the buttons multiply: every
component of the fare stack, from the metered fare the driver keeps to
the congestion fees the driver merely collects and remits, lands in one
total, and the percentages are computed on the sum. This is not one
vendor's quirk. All three payment vendors match tips to the
fee-inclusive base at four to five times the rate of any narrower base
(Table~\ref{tab:vendor}), and the base is visible in behavior alone:
57 percent of positive credit-card tips sit in the half-point bin at
\emph{exactly} 20 percent of the all-in total, and the same tips,
recomputed against the metered fare, show no spikes at all
(Figure~\ref{fig:spikes}). \textcite{haggag2014default} established
that suggested tips cause tips. A suggestion, however, is a percentage
\emph{of something}. This paper is about the something---specifically,
about what happens when public policy moves it.

The question is the pass-through rate: when a statute puts one more
dollar on the meter, how many cents cross into the tip? Call it $\rho$.
Answering requires surcharge variation that is not correlated with who
is riding, and New York's fare schedule supplies it daily: the
rush-hour surcharge applies to weekday pickups from 4:00pm, assigned by
the pickup minute. We build date-by-minute cells from the universe of
standard-rate credit-card trips, September 2021 through September 2024,
and compare tips minutes apart across the cutoffs.

\subsection{A daily experiment, and a trap in it}
Tips jump at 4:00pm. But the naive comparison is a trap: the metered
fare jumps too, by \$0.49, because rush-hour traffic makes trips slower
and dearer (Table~\ref{tab:balance}), and attributing the whole tip step
to the surcharge inflates it by a third (\$0.389 against \$0.292).
The repair is in the fare schedule itself. On weekends the same clock
carries no surcharge, so the weekday--weekend \emph{difference in
discontinuities} at 4:00pm nets out whatever rush hour does to trips on
any day of the week \parencite{grembi2016difference}; our analysis plan
locked weekends as controls before estimation. Conditioning on the
non-surcharge fare, the weekend placebo step falls from \$0.088 to
\$0.024; differencing weekday against weekend gives $0.108$ per
statutory dollar in the \$2.50 regime and $0.105$ in the \$1.00 regime,
and dropping the fare control entirely moves those by about a cent.
Identification does not lean on a functional-form assumption about
fares; the weekend difference does the work.

\subsection{One parameter}
The headline result is that about \emph{eleven cents of every statutory
surcharge dollar becomes tip}. Across the six primary within-day
changes---the rush surcharge switching on at two statutory doses, the
overnight surcharge switching on at two doses, and two downward changes
where the net surcharge falls---the intention-to-treat effect per
statutory dollar spans $0.090$ to $0.117$, with a precision-weighted
mean of $0.108$ (SE $0.001$). Scaled instead by the surcharge dollars
the meter actually records, $\rho^{IV}$ runs $0.122$ to $0.150$. The
average tip in this sample is 17.4 percent of the base, so riders tip a
surcharge dollar at roughly four-fifths the rate they tip a fare
dollar: the surcharge is not distinguished from the fare because, on
the screen, it is not distinguishable from the fare.

Figure~\ref{fig:coefplot} assembles every locked estimate on one axis,
and it is the paper's argument in one frame. Two statutory doses,
because the legislature repriced both surcharges by 150 percent in
December 2022. Surcharges switching on and switching off. Three clocks.
Weekday and weekend cutoffs. A second control series---legal holidays,
working days with no surcharge---replicates the weekend difference at
$0.095$--$0.096$, and its exception proves the rule: Juneteenth, which
the meters do not exempt, shows the full step. A \$5.00 surcharge on
JFK \emph{flat} fares---where the metered fare cannot jump at 4:00pm by
construction---gives $0.133$--$0.149$ per statutory dollar with no
controls at all, twice our largest dose on a base three times larger.
And---with no clock in it at all---the January 2025 congestion toll,
identified from a boundary in space: the spatial estimate is $0.107$
per statutory dollar, against the clock designs' pooled $0.108$. Two estimates fail: both at 6:00am on
weekdays, where the dawn shift toward airport-bound trips breaks the
design's conditioning. We report them in grey rather than discard them,
because their weekend twins at the same clock and dose are
normal---which localizes the failure to weekday-dawn composition, not
to the mechanism, and shows the identifying assumption is checkable
(Section~\ref{sec:results-6am}).

\subsection{Riders do not undo it}
Three findings say the pass-through is a property of the interface
rather than of attentive generosity. First, the button is not the
mechanism: the two major vendors' riders take an exact default on 19
versus 65 percent of tips, yet pass-through is statistically identical
at both ($0.147$ and $0.152$ per recorded dollar)---riders who type
custom amounts compute percentages on the same inclusive total the
screen displays. Second, riders barely react to repricing: when the
rush surcharge rose from \$1.00 to \$2.50 overnight in December 2022,
keeping tip dollars fixed required the rush-window tip rate to fall
1.1 points; it fell 0.22, immediately, and never drifted further over
the following six months. Third, the quantity margins are silent:
pickup density is smooth through the cutoff, card use does not fall,
and zero-tipping rises 0.15 percentage points. Riders absorb the fee on
every margin except the one the software moves for them.

\subsection{What is new}
That payment interfaces differ in the base they multiply, and that
add-ons inside the base raise tips, is established:
\textcite{thakral2026five} document both in this market, using the 2012
fare change and the two archival vendors' differing treatment of tax
and tolls, and our archival evidence confirms their account---the base
difference explains about three-quarters of the vendor tip gap that
\textcite{haggag2014default} used to identify default effects
(Appendix~\ref{app:replication}). What has not been measured is what
\emph{statutory policy} does when it rides that base. A legislated
surcharge that switches at a known minute, at a dose the legislature
set and then changed, is variation of a different kind than a fare
change or a route-determined toll, and it identifies a policy object:
the pass-through of a public fee into a private gratuity, per statutory
dollar. We estimate it, test its stability across doses, directions,
clocks, and two congestion-pricing programs, and aggregate the
resulting transfer. Prior work shows that interfaces construct the
tipping norm; this paper shows that fiscal policy conscripts it.

\subsection{Incidence, stakes, and a policy dial}
The driver remits the congestion fee and keeps the tip, so the toll's
true incidence splits: on a \$0.75 CBD toll, the rider pays about
\$0.83, the MTA receives \$0.75, and the driver nets the induced tip.
Aggregated against fees actually charged, the 2019 congestion
surcharge's \$127.2 million on card trips generated roughly
\$18.7 million a year in unlegislated gratuities, and the 2025 toll
adds about \$2.2 million a year (Section~\ref{sec:counterfactuals}).
These are transfers, not losses---and they run from riders to drivers,
plausibly progressively. The point is not that the transfer is bad; it
is that no one chose it. The dial that sets it is one line of
payment-terminal arithmetic---the definition of the base---and the 2012
market, in which one vendor computed on a narrower base, shows the dial
is real engineering, not hypothesis. What a default is computed
\emph{on} matters as much as what it is set \emph{to}.

The rest of the paper proceeds as follows.
Section~\ref{sec:related-literature} situates the paper.
Section~\ref{sec:empirical-facts} presents the fare stack, the data,
and the base facts. Section~\ref{sec:model} states the designs.
Section~\ref{sec:results} reports the estimates, the mechanism
evidence, and the congestion-fee designs.
Section~\ref{sec:counterfactuals} prices the transfers and the
counterfactuals, and Section~\ref{sec:conclusion} concludes.
